\documentclass[11pt,a4paper]{article}
\usepackage{fontspec}
\setmainfont{DejaVu Sans}[Scale=0.82]
\usepackage{amsmath,amssymb}
\usepackage[left=16mm,right=16mm,top=7mm,bottom=6mm]{geometry}
\usepackage[table]{xcolor}
\usepackage{array,tabularx,booktabs}
\definecolor{grau}{gray}{0.45}
\definecolor{band}{gray}{0.94}
\pagestyle{empty}
\setlength{\parindent}{0pt}
\renewcommand{\arraystretch}{1.0}
\setlength{\tabcolsep}{7pt}
\newcommand{\g}[1]{{\small\color{grau}#1}}
\newcommand{\block}[1]{\hiderowcolors\multicolumn{2}{@{}l}{\rule{0pt}{9pt}\color{grau}\small\bfseries\MakeUppercase{#1}}\\[-2pt]\showrowcolors}
\newcolumntype{S}{>{\bfseries\raggedright\arraybackslash}p{5.0cm}}
\newcolumntype{G}{>{\raggedright\arraybackslash}p{5.3cm}}
\newcommand{\grenz}[1]{{\footnotesize\color{grau}#1}}
\newcolumntype{A}{>{\raggedright\arraybackslash}X}
\begin{document}
{\Large\bfseries Kurvendiskussion}\hfill{\color{grau}Stichwort → Ansatz \quad Grenze → Ersatz}\par
\vspace{3pt}\hrule height 1pt\vspace{6pt}
\g{$x_N$ Nullstelle · $x_E$ Extremstelle · $x_W$ Wendestelle · $x_S$ Sattelstelle.
Stelle = $x$-Wert, Punkt = $(x\,|\,y)$.}
\vspace{4pt}

\rowcolors{1}{band}{white}
\begin{tabularx}{\linewidth}{@{}SA@{}}
\block{Bereiche und Symmetrie}
Definitionsbereich & verboten: Nenner $=0$, Radikand $<0$, Argument von $\ln\le0$\newline\grenz{Grenze: ganzrational: $D=\mathbb R$, nichts zu prüfen}\\
Wertebereich & vergleiche: Grenzwerte am Rand von $D$, \ $f(x_E)$, \ Randwerte\newline\grenz{Grenze: nur mit allen Extrempunkten → sonst aus Graph und Monotonie}\\
Symmetrie zur $y$-Achse & $f(x)=f(-x)$ \quad \g{ganzrational: nur gerade Exponenten}\newline\grenz{Grenze: Exponenten-Regel nur ganzrational → sonst $f(-x)$ einsetzen}\\
Symmetrie zum Ursprung & $f(x)=-f(-x)$ \quad \g{ganzrational: nur ungerade Exponenten}\newline\grenz{Grenze: Exponenten-Regel nur ganzrational → sonst $-f(-x)$ einsetzen}\\
keine Symmetrie & weder noch \quad \g{gilt nur für $y$-Achse und $O$}\newline\grenz{Grenze: nachweisen mit Gegenbeispiel: $f(1)\neq\pm f(-1)$}\\
\block{Schnittpunkte und Unendlich}
Punktprobe & $y_P\overset{?}{=}f(x_P)$\\
Nullstellen & $f(x_N)=0$\newline\grenz{Grenze: Grad $\ge3$: ausklammern, raten + Division → sonst $f(a)>0$, $f(b)<0$ ⇒ $x_N$ dazwischen}\\
Schnittpunkte mit den Achsen & $S_x(x_N\,|\,0)$, \quad $S_y(0\,|\,f(0))$\\
Schnittpunkte zweier Graphen & $f(x)=g(x)$, \ dann $y=f(x)$\newline\grenz{Grenze: wie Nullstellen von $f-g$}\\
Verhalten im Unendlichen & $\lim\limits_{x\to\pm\infty}f(x)$: \ höchste Potenz entscheidet \quad \g{$\mathrm e^x$ wächst schneller als jede Potenz}\newline\grenz{Grenze: „höchste Potenz“ nur ganzrational → $\mathrm e^x$-Regel}\\
\block{Ableitung}
Anstieg an der Stelle $x_0$ & $m=f'(x_0)$\\
Steigungswinkel & $\tan\alpha=f'(x_0)$\newline\grenz{Grenze: $m<0$: $\alpha=180^\circ+\arctan m$}\\
mittlere Änderungsrate & $\dfrac{f(b)-f(a)}{b-a}$ \quad \g{Sekantensteigung}\\
momentane Änderungsrate & $f'(x_0)$ \quad \g{Einheit: Einheit von $f$ je Einheit von $x$}\\
Extremstellen & $f'(x_E)=0$ \ und \ $f''(x_E)\neq0$\newline\grenz{Grenze: $f''(x_E)=0$: kein Schluss → VZW von $f'$}\\
Hochpunkt & $f'(x_E)=0$ \ und \ $f''(x_E)<0$: \ $H(x_E\,|\,f(x_E))$\\
Tiefpunkt & $f'(x_E)=0$ \ und \ $f''(x_E)>0$: \ $T(x_E\,|\,f(x_E))$\\
Wendestellen & $f''(x_W)=0$ \ und \ $f'''(x_W)\neq0$\newline\grenz{Grenze: $f'''(x_W)=0$: kein Schluss → VZW von $f''$}\\
Wendepunkte & wie Wendestellen: \ $W(x_W\,|\,f(x_W))$\\
Sattelstellen & $f'(x_S)=0$, \ $f''(x_S)=0$, \ $f'''(x_S)\neq0$\newline\grenz{Grenze: $f'''(x_S)=0$: kein Schluss → $f'$ ohne VZW ⇒ Sattel}\\
Sattelpunkt & wie Sattelstellen: \ $S(x_S\,|\,f(x_S))$\\
Vorzeichenwechsel & statt $f''$ bzw.\ $f'''$: \ $f'$: $+\to-$ $H$, \ $-\to+$ $T$, \ kein Wechsel $S$\newline\grenz{Grenze: Testwert zwischen je zwei Nullstellen von $f'$}\\
\block{Verlauf}
Monotonie & $f'>0$ steigt, \ $f'<0$ fällt \quad \g{Grenzen: Extremstellen}\newline\grenz{Grenze: Sattelstelle trennt nicht: Bereiche zusammenfassen}\\
Krümmung & $f''>0$ links, \ $f''<0$ rechts \quad \g{Grenzen: Wendestellen}\\
globale Extrema & $f(x_E)$ mit Randwerten und Grenzwerten vergleichen\newline\grenz{Grenze: Grenzwert $\pm\infty$ ⇒ kein globales Extremum}\\
stärkste Zu- oder Abnahme & im Wendepunkt (Extremum von $f'$)\newline\grenz{Grenze: im Intervall auch $f'$ an den Rändern}\\
Graph von $f$ aus Graph von $f'$ & Nullstelle von $f'$ mit VZW: Extremstelle von $f$; \ Extremstelle von $f'$: Wendestelle von $f$\newline\grenz{Grenze: Nullstelle von $f'$ ohne VZW: Sattelstelle}\\
Tangente an der Stelle $x_0$ & $t(x)=mx+n$, \ $m=f'(x_0)$, \ $n$ aus $t(x_0)=f(x_0)$\\
Normale an der Stelle $x_0$ & $n(x)=m_nx+b$, \ $m_n=-\dfrac{1}{f'(x_0)}$, \ $b$ aus $n(x_0)=f(x_0)$\newline\grenz{Grenze: $f'(x_0)=0$: Normale $x=x_0$}\\
\end{tabularx}
\end{document}
